% Section 1: Introduction

\section{Introduction}
\label{sec:introduction}

\lettrine[lines=2, lhang=0.15, loversize=0.1]{T}{he defining feature} of modern investment funds is not their asset strategy, but their documentary complexity. A standard Luxembourg Reserved Alternative Investment Fund (\raif{}) structure is not a single entity but a \textit{network} of interlocking legal agreements: a Limited Partnership Agreement (LPA), an AIFM Agreement, a Depositary Agreement, an Administration Agreement, and dozens of side letters granting individual investors bespoke terms. Each document imposes constraints on the others; a fee cap negotiated in a side letter must propagate to the Net Asset Value calculation logic in the Administration Agreement, which must align with the waterfall provisions in the LPA.

Currently, the consistency of this distributed system is verified manually. This approach is fundamentally non-scalable. As the number of documents ($N$) grows, the number of potential cross-references grows as \bigO{N^2}. For a typical structure with 50+ documents, manual verification relies on heuristics---trust, precedent, and sample-checking---rather than exhaustive proof. This creates what we term \textbf{``legal bugs''}: silent inconsistencies that remain dormant until a liquidity event, regulatory audit, or investor dispute triggers a crisis.

\subsection{The Case for Formal Verification}

We posit that fund law is amenable to \textbf{formal verification}. Unlike criminal law, which relies heavily on reasonable person standards, or constitutional law, which involves balancing abstract principles, fund law is largely \textit{structural} and \textit{arithmetic}. It deals with:

\begin{itemize}[nosep,leftmargin=1.5em]
  \item \textbf{Thresholds}: Assets under Management $\geq$ €100M triggers full \aifmd{} scope
  \item \textbf{Binary states}: Authorized vs.\ Registered AIFM status
  \item \textbf{Precise definitions}: Professional Investor per \mifid{} Annex~II
  \item \textbf{Algorithmic calculations}: Waterfall distributions, carried interest
\end{itemize}

These characteristics suggest that significant portions of fund law can be encoded as executable specifications that compile to provably correct implementations.

\subsection{Contributions}

In this work, we present \textbf{Sanctaphrax}, a proof-of-concept formalization of Luxembourg fund law using \textbf{\catala{}}~\cite{merigoux2021catala}. \catala{} is a domain-specific language (DSL) for literate programming that compiles law into correct-by-construction executable artifacts. It uses \textbf{default logic} to handle the general-case/exception structure typical of legislative texts---precisely the structure that dominates fund regulation.

Our contributions are as follows:

\begin{enumerate}[nosep,leftmargin=1.5em]
  \item We provide the first formal specification of the \textbf{Definition of an AIF} under \aifmd{} Article~4(1)(a), explicitly documenting the interpretive choices required.

  \item We implement a \textbf{Professional Investor eligibility oracle} that unifies \mifid{} and \raif{} Law requirements into a single decidable procedure.

  \item We demonstrate an automated \textbf{Cross-Reference Verification} mechanism that detects definition inconsistencies across separate fund documents.

  \item We analyze the \textbf{limits of formalization}, identifying and categorizing the approximately 15\% of provisions where human judgment cannot be replaced by code.
\end{enumerate}

\subsection{Approach and Limitations}

Unlike Large Language Model (LLM) approaches which \textit{predict} likely legal text based on statistical patterns, our approach provides \textbf{provable guarantees}. If the specification compiles and the inputs match the real-world situation, the conclusion follows \textit{deductively} from the encoded law. There is no hallucination, no probabilistic uncertainty in the reasoning chain.

However, we are explicit about what formal verification does \textit{not} provide:

\begin{itemize}[nosep,leftmargin=1.5em]
  \item It does not guarantee that our encoding matches the ``true'' legal meaning
  \item It does not replace the judgment required for open-textured concepts
  \item It does not predict how a court or regulator would interpret ambiguous provisions
\end{itemize}

The value lies in \textit{locating} where human judgment is required, not in replacing it. A specification that flags ``these two definitions of Permitted Investment differ'' allows a human to determine whether the difference is acceptable. The machine surfaces the question; the lawyer answers it.

\subsection{Paper Organization}

Section~\ref{sec:background} introduces \catala{} and default logic, explaining why fund law is a suitable domain. Sections~\ref{sec:aif}--\ref{sec:crossref} present our four specifications: AIF definition, Professional Investor eligibility, \raif{} structural requirements, and cross-reference verification. Section~\ref{sec:findings} analyzes where formalization breaks down. Section~\ref{sec:implications} discusses implications for practitioners, regulators, and researchers. Section~\ref{sec:questions} poses open questions. Section~\ref{sec:conclusion} concludes.
